% Part: first-order-logic
% Chapter: syntax-and-semantics
% Section: assignments

\documentclass[../../../include/open-logic-section]{subfiles}

\begin{document}

\olfileid{fol}{syn}{ass}

\olsection{Toekenningen aan variabelen}

\begin{explain}
Een toekenning aan variabelen~$s$ geeft \emph{elke} variabele een
waarde. Er zijn oneindig veel variabelen. Voor één bepaalde term of
formule hebben we lang niet al die waarden nodig. Toch geven we elke
variabele een waarde, omdat de regels daardoor eenvoudiger worden.
De waarde van een term~$t$ en de vraag of !!a{formula}~$!A$ waar is in
!!a{structure} bij~$s$, hangen alleen af van wat de toekenning~$s$
geeft aan de !!{variable}s in~$t$ en aan de vrije !!{variable}s
van~$!A$. De volgende twee uitspraken maken dat precies. Als twee
toekenningen hetzelfde geven aan alle variabelen in~$t$, komt de term
bij beide op dezelfde waarde uit. $!A$ is bij de ene toekenning precies
dan waar als bij de andere wanneer ze hetzelfde geven aan alle vrije
variabelen van~$!A$.
\end{explain}

\begin{prop}\ollabel{prop:valindep}
Staan in een term~$t$ alleen de !!{variable}s $x_1$, \dots,~$x_n$,
en geldt $s_1(x_i) = s_2(x_i)$ voor $i = 1$, \dots,~$n$, dan geldt
$\Value{t}{M}[s_1] = \Value{t}{M}[s_2]$.
\end{prop}

\begin{proof}
We gebruiken inductie over de opbouw van~$t$. Begin met de eenvoudigste
termen. Dan is~$t$ !!a{constant} of een van de variabelen~$x_1$,
\dots,~$x_n$. Als $t = c$, dan
$\Value{t}{M}[s_1] = \Assign{c}{M} = \Value{t}{M}[s_2]$. Als
$t = x_i$, dan weten we al dat $s_1(x_i) = s_2(x_i)$. Daarom
$\Value{t}{M}[s_1] = s_1(x_i) = s_2(x_i) = \Value{t}{M}[s_2]$.

Neem nu een langere term $t = \Atom{f}{t_1, \dots, t_k}$. Neem aan
dat de uitspraak al klopt voor de onderdelen $t_1$, \dots, $t_k$.
Dit is de inductiehypothese. Dan
\begin{align*}
  \Value{t}{M}[s_1] & = \Value{\Atom{f}{t_1, \dots, t_k}}{M}[s_1] \\
  & = \Assign{f}{M}(\Value{t_1}{M}[s_1], \dots, \Value{t_k}{M}[s_1]).
\intertext{Elk onderdeel~$t_j$, met $j = 1$, \dots,~$k$, bevat alleen
    !!{variable}s uit $x_1$, \dots,~$x_n$. De inductiehypothese geeft
    dus $\Value{t_j}{M}[s_1] = \Value{t_j}{M}[s_2]$. Daarom}
\Value{t}{M}[s_1] & = \Value{\Atom{f}{t_1, \dots, t_k}}{M}[s_1] \\
  & = \Assign{f}{M}(\Value{t_1}{M}[s_1], \dots, \Value{t_k}{M}[s_1]) \\
  & = \Assign{f}{M}(\Value{t_1}{M}[s_2], \dots, \Value{t_k}{M}[s_2]) \\
  & = \Value{\Atom{f}{t_1, \dots, t_k}}{M}[s_2] = \Value{t}{M}[s_2].
\end{align*}
\end{proof}

\begin{prop}\ollabel{prop:satindep}
Zijn de vrije !!{variable}s in~$!A$ allemaal een van
$x_1$, \dots,~$x_n$, en geldt $s_1(x_i) = s_2(x_i)$ voor
$i = 1$, \dots,~$n$, dan geldt $\Sat{M}{!A}[s_1]$ precies als
$\Sat{M}{!A}[s_2]$.
\end{prop}

\begin{proof}
Ook hier gebruiken we inductie over de opbouw, nu van~$!A$. Begin met
een atomaire formule. Dan is~$!A$ atomair en kan~$!A$ een van deze
vormen hebben:
\iftag{prvTrue}{$\ltrue$,}{}
\iftag{prvFalse}{$\lfalse$,}{}
$\Atom{R}{t_1, \dots, t_k}$ met een $k$-plaatsig predicaat~$R$ en de
termen $t_1$, \dots,~$t_k$, of $\eq[t_1][t_2]$ met de termen $t_1$
en~$t_2$. Bij de laatste twee gevallen bewijzen we alleen de ene
richting van de !!{biconditional}. De andere richting werkt op
dezelfde manier.

\begin{enumerate}
\tagitem{prvTrue}{%
  \indcase{!A}{\ltrue}{bij beide toekenningen geldt
    $\Sat{M}{\indfrm}[s_1]$ en $\Sat{M}{\indfrm}[s_2]$.}}{}

\tagitem{prvFalse}{%
  \indcase{!A}{\lfalse}{bij beide toekenningen geldt
    $\Sat/{M}{!A}[s_1]$ en $\Sat/{M}{!A}[s_2]$.}}{}

\item
  \indcase{!A}{\Atom{R}{t_1, \ldots, t_k}}{neem aan dat
    $\Sat{M}{\indfrm}[s_1]$. Dan
    \[
    \langle \Value{t_1}{M}[s_1], \ldots, \Value{t_k}{M}[s_1] \rangle
    \in \Assign{R}{M}.
    \]
    Voor elk $i = 1$, \dots,~$k$ zegt \olref{prop:valindep} dat
    $\Value{t_i}{M}[s_1] = \Value{t_i}{M}[s_2]$. We kunnen dus elk
    onderdeel van de tuple vervangen en krijgen
    $\langle \Value{t_1}{M}[s_2], \ldots, \Value{t_k}{M}[s_2] \rangle
    \in \Assign{R}{M}$. Daarom $\Sat{M}{\indfrm}[s_2]$.}
\item
  \indcase{!A}{\eq[t_1][t_2]}{neem aan dat
    $\Sat{M}{\indfrm}[s_1]$. Dan is
    $\Value{t_1}{M}[s_1] = \Value{t_2}{M}[s_1]$. Nu volgt stap voor stap:
    \begin{align*}
      \Value{t_1}{M}[s_2] & = \Value{t_1}{M}[s_1]
      & \text{(volgens \olref{prop:valindep})} \\
      & = \Value{t_2}{M}[s_1] 
      & \text{(omdat $\Sat{M}{\eq[t_1][t_2]}[s_1]$)}\\
      &= \Value{t_2}{M}[s_2]
      & \text{(volgens \olref{prop:valindep}).}
    \end{align*}
    Dus $\Sat{M}{\eq[t_1][t_2]}[s_2]$.}
\end{enumerate}

Neem nu aan dat $\Sat{M}{!B}[s_1]$ precies dan geldt als
$\Sat{M}{!B}[s_2]$ voor elke eenvoudiger !!{formula}~$!B$ dan~$!A$.
Dat is de inductiehypothese. We kijken naar het buitenste logische
teken van~$!A$, de hoofdoperator. Bij elk geval bewijzen we alleen de
ene richting; terug werkt het op dezelfde manier. Behalve bij de
kwantoren gebruiken we de inductiehypothese voor een kleinere
!!{formula}~$!B$ in~$!A$. Elke vrije variabele van~$!B$ is ook vrij
in~$!A$. Als $s_1$ en $s_2$ hetzelfde geven aan de vrije variabelen
van~$!A$, doen ze dat dus ook aan die van~$!B$, en mogen we de
inductiehypothese op~$!B$ toepassen.

\begin{enumerate}
\tagitem{defNot}{}{%
  \iftag{probNot}{%
    \indcase!{!A}{\lnot !B}{}}{%
    \indcase{!A}{\lnot !B}{neem aan dat
      $\Sat{M}{\indfrm}[s_1]$. Dan is de formule bij die toekenning
      niet waar: $\Sat/{M}{!B}[s_1]$. Volgens de inductiehypothese is
      zij ook bij de andere toekenning niet waar:
      $\Sat/{M}{!B}[s_2]$. Daarom
      $\Sat{M}{\indfrm}[s_2]$.}}}

\tagitem{defAnd}{}{%
  \iftag{probAnd}{%
    \indcase!{!A}{!B \land !C}{}}{%
    \indcase{!A}{!B \land !C}{neem aan dat
      $\Sat{M}{\indfrm}[s_1]$. Dan gelden
      $\Sat{M}{!B}[s_1]$ en $\Sat{M}{!C}[s_1]$. De
      inductiehypothese geeft
      $\Sat{M}{!B}[s_2]$ en $\Sat{M}{!C}[s_2]$. Daarom
      $\Sat{M}{\indfrm}[s_2]$.}}}

\tagitem{defOr}{}{%
  \iftag{probOr}{%
    \indcase!{!A}{!B \lor !C}{}}{%
    \indcase{!A}{!B \lor !C}{neem aan dat
      $\Sat{M}{\indfrm}[s_1]$. Dan geldt $\Sat{M}{!B}[s_1]$ of
      $\Sat{M}{!C}[s_1]$. De inductiehypothese geeft
      $\Sat{M}{!B}[s_2]$ of
      $\Sat{M}{!C}[s_2]$. Dus $\Sat{M}{\indfrm}[s_2]$.}}}

\tagitem{defIf}{}{%
  \iftag{probIf}{%
    \indcase!{!A}{!B \lif !C}{}}{%
    \indcase{!A}{!B \lif !C}{neem aan dat
    $\Sat{M}{\indfrm}[s_1]$. Dan is het linkerdeel niet waar of het
    rechterdeel wel waar: $\Sat/{M}{!B}[s_1]$ of
    $\Sat{M}{!C}[s_1]$. De inductiehypothese geeft
    $\Sat/{M}{!B}[s_2]$ of $\Sat{M}{!C}[s_2]$. Daarom
    $\Sat{M}{\indfrm}[s_2]$.}}}

\tagitem{defIff}{}{%
  \iftag{probIff}{%
    \indcase!{!A}{!B \liff !C}{}}{%
    \indcase{!A}{!B \liff !C}{neem aan dat
      $\Sat{M}{\indfrm}[s_1]$. Dan zijn beide delen allebei waar of
      allebei niet waar. In symbolen: óf
      $\Sat{M}{!B}[s_1]$ en $\Sat{M}{!C}[s_1]$, óf
      $\Sat/{M}{!B}[s_1]$ en $\Sat/{M}{!C}[s_1]$. De
      inductiehypothese geeft bij de andere toekenning hetzelfde
      patroon: óf
      $\Sat{M}{!B}[s_2]$ en $\Sat{M}{!C}[s_2]$, óf
      $\Sat/{M}{!B}[s_2]$ en $\Sat/{M}{!C}[s_2]$. In beide gevallen
      geldt $\Sat{M}{\indfrm}[s_2]$.}}}

\tagitem{defEx}{}{%
  \iftag{probEx}{%
    \indcase!{!A}{\lexists[x][!B]}{}}{%
    \indcase{!A}{\lexists[x][!B]}{neem aan dat
      $\Sat{M}{\indfrm}[s_1]$. Dan is er minstens één
      $m \in \Domain{M}$ waarvoor
      $\Sat{M}{!B}[\Subst{s_1}{m}{x}]$. Noem de aangepaste
      toekenningen $s_1' = \Subst{s_1}{m}{x}$ en $s_2' =
      \Subst{s_2}{m}{x}$. De vrije variabelen van~$!B$ zitten tussen
      $x_1$, \dots, $x_n$ en~$x$. Er geldt
      $s_1'(x_i) = s_2'(x_i)$, want $s_1'$ en $s_2'$ zijn
      $x$-varianten van $s_1$ en~$s_2$: ze kunnen alleen bij de
      genoemde variabele anders zijn. Volgens de aanname is
      $s_1(x_i) = s_2(x_i)$. Ook geven de twee aangepaste
      toekenningen dezelfde nieuwe waarde:
      $s_1'(x) = s_2'(x) = m$. Dat volgt uit de manier waarop
      $s_1'$ en~$s_2'$ zijn gemaakt. De inductiehypothese geldt dus
      voor $!B$ bij $s_1'$ en~$s_2'$, en geeft
      $\Sat{M}{!B}[s_2']$. Omdat $s_2' = \Subst{s_2}{m}{x}$, is er
      een $m \in \Domain{M}$ waarvoor
      $\Sat{M}{!B}[\Subst{s_2}{m}{x}]$. Daarom
      $\Sat{M}{\indfrm}[s_2]$.}}}

\tagitem{defAll}{}{%
  \iftag{probAll}{%
    \indcase!{!A}{\lforall[x][!B]}{}}{%
    \indcase{!A}{\lforall[x][!B]}{neem aan dat
      $\Sat{M}{\indfrm}[s_1]$. Dan geldt voor elke
      $m \in \Domain{M}$ dat
      $\Sat{M}{!B}[\Subst{s_1}{m}{x}]$. We moeten bewijzen dat ook
      voor elke $m \in \Domain{M}$ geldt:
      $\Sat{M}{!B}[\Subst{s_2}{m}{x}]$. Kies nu een willekeurige
      $m \in \Domain{M}$ en maak de twee aangepaste toekenningen
      $s_1' = \Subst{s_1}{m}{x}$ en $s_2' =
      \Subst{s_2}{m}{x}$. Voor de eerste geldt
      $\Sat{M}{!B}[s_1']$. De vrije variabelen van~$!B$ zitten tussen
      $x_1$, \dots, $x_n$ en~$x$. Er geldt
      $s_1'(x_i) = s_2'(x_i)$, want $s_1'$ en $s_2'$ zijn
      $x$-varianten van $s_1$ en~$s_2$. Volgens de aanname is
      $s_1(x_i) = s_2(x_i)$. Bij de genoemde variabele geven ze
      allebei dezelfde waarde: $s_1'(x) = s_2'(x) = m$. De aldus
      gemaakte toekenningen $s_1'$ en~$s_2'$ voldoen dus aan de
      voorwaarden. De inductiehypothese op~$!B$ met~$s_1'$ en~$s_2'$
      geeft $\Sat{M}{!B}[s_2']$. Dit geldt voor elke
      $m \in \Domain{M}$; dat wil zeggen:
      $\Sat{M}{!B}[\Subst{s_2}{m}{x}]$ voor alle~$m \in
      \Domain{M}$. Dus $\Sat{M}{\indfrm}[s_2]$.}}}
\end{enumerate}
De inductie is daarmee klaar. Dus geldt
$\Sat{M}{!A}[s_1]$ precies als $\Sat{M}{!A}[s_2]$ wanneer de vrije
!!{variable}s in~$!A$ allemaal een van $x_1$, \dots, $x_n$ zijn en
$s_1(x_i)=s_2(x_i)$ voor $i = 1$, \dots,~$n$.
\end{proof}

\begin{probtag}{probNot,probOr,probAnd,probIf,probIff,probEx,probAll}
Maak het bewijs van \olref[fol][syn][ass]{prop:satindep} af.
\end{probtag}

\begin{explain}
Een !!{sentence} is een formule zonder vrije variabelen. Er zijn dus
nul vrije variabelen waarover twee toekenningen het oneens kunnen
zijn. Elke twee toekenningen stemmen daar automatisch op overeen. De
vorige uitspraak zegt daarom iets eenvoudigs: of een zin waar is in
een structuur, hangt helemaal niet af van de gekozen toekenning. We
leggen dat hieronder vast. Daarna kunnen we zeggen dat !!a{sentence}
waar is in !!a{structure} zonder er steeds een toekenning bij te noemen.
\end{explain}

\begin{cor}
\ollabel{cor:sat-sentence}
Als $!A$ !!a{sentence} is en $s$ een toekenning aan variabelen, dan
geldt $\Sat{M}{!A}[s]$ precies als $\Sat{M}{!A}[s']$ voor elke andere
toekenning~$s'$.
\end{cor}

\begin{proof}
  Kies een willekeurige toekenning~$s'$. Omdat $!A$ een zin is, heeft
  zij geen vrije variabelen. Elke toekenning~$s'$ geeft dus vanzelf
  hetzelfde aan alle vrije variabelen van~$!A$ als~$s$: die zijn er
  niet. Daarmee is aan de voorwaarde van
  \olref{prop:satindep} voldaan. Dus geldt $\Sat{M}{!A}[s]$ precies als
  $\Sat{M}{!A}[s']$.
\end{proof}

\begin{defn}
\ollabel{defn:satisfaction}
Als $!A$ !!a{sentence} is, zeggen we dat !!a{structure}~$\Struct M$
$!A$ \emph{waar maakt}. We schrijven $\Sat{M}{!A}$. Dit betekent dat
$\Sat{M}{!A}[s]$ geldt voor elke toekenning aan variabelen~$s$.
\end{defn}

Als $\Sat{M}{!A}$, zeggen we ook kortweg dat \emph{$!A$ waar is
in~$\Struct{M}$.} We kunnen hetzelfde idee gebruiken voor een hele
verzameling !!{sentence}s.

\begin{defn}
\ollabel{defn:sat}
Neem een verzameling~$\Gamma$ van !!{sentence}s; dezelfde verzameling
noemen we verder~$\Gamma$. We zeggen dat !!a{structure}~$\Struct M$
de verzameling~$\Gamma$ \emph{waar maakt},
en schrijven $\Sat{M}{\Gamma}$, als $\Sat{M}{!A}$ geldt voor elke
$!A \in \Gamma$.
\end{defn}

\begin{prop}\ollabel{prop:sentence-sat-true}
  Neem !!a{structure}~$\Struct{M}$, !!a{sentence}~$!A$ en een
  toekenning aan variabelen~$s$. Dan geldt $\Sat{M}{!A}$ precies als
  $\Sat{M}{!A}[s]$.
\end{prop}

\begin{proof}
Opgave.
\end{proof}

\begin{prob}
Bewijs \olref[fol][syn][ass]{prop:sentence-sat-true}.
\end{prob}

\begin{prop}\ollabel{prop:sat-quant}
  Stel dat in~$!A(x)$ alleen $x$ vrij voorkomt, en neem
  !!a{structure}~$\Struct M$. Dan:
  \begin{tagenumerate}{prvEx,prvAll}
    \tagitem{prvEx}{$\Sat{M}{\lexists[x][!A(x)]}$ precies als
      $\Sat{M}{!A(x)}[s]$ voor minstens één toekenning~$s$.}{}
    \tagitem{prvAll}{$\Sat{M}{\lforall[x][!A(x)]}$ precies als
      $\Sat{M}{!A(x)}[s]$ voor elke toekenning~$s$.}{}
\end{tagenumerate}
\end{prop}

\begin{proof}
Opgave.
\end{proof}

\begin{prob}
Bewijs \olref[fol][syn][ass]{prop:sat-quant}.
\end{prob}

\begin{prob}
\DeclareRobustCommand{\VDash}{\mathrel{||}\joinrel\Relbar}  
Stel dat de taal~$\Lang L$ geen !!{function}s heeft. Neem een
!!{structure}~$\Struct M$, !!a{constant}~$c$ en een object
$a \in \Domain M$. Met $\Struct M[a/c]$ bedoelen we een kopie van
$\Struct M$ waarin alleen de betekenis van die constante is veranderd:
$\Assign{c}{M[a/c]} = a$. Voor $\Struct M \VDash !A$, waarbij
!!a{sentence}~$!A$ is, gelden nu rechtstreeks de volgende regels:
\begin{enumerate}
\tagitem{prvFalse}{%
  \indcase{!A}{\lfalse}{niet $\Struct{M} \VDash \indfrm$.}}{}

\tagitem{prvTrue}{%
  \indcase{!A}{\ltrue}{$\Struct{M} \VDash \indfrm$.}}{}

\item \indcase{!A}{\Atom{R}{d_1, \dots, d_n}}{$\Struct M \VDash \indfrm$
  precies als de bijbehorende rij objecten in de juiste relatie staat:
  $\langle \Assign{d_1}{M}, \dots, \Assign{d_n}{M} \rangle \in
  \Assign{R}{M}$.}

\item \indcase{!A}{\eq[d_1][d_2]}{$\Struct M \VDash \indfrm$ precies
  als $\Assign{d_1}{M} = \Assign{d_2}{M}$.}

\tagitem{prvNot}{%
  \indcase{!A}{\lnot !B}{$\Struct{M} \VDash \indfrm$ precies als niet
    $\Struct{M} \VDash {!B}$.}}{}

\tagitem{prvAnd}{%
  \indcase{!A}{(!B \land !C)}{$\Struct{M} \VDash
    \indfrm$ precies als $\Struct{M} \VDash!B$ en
    $\Struct{M} \VDash !C$.}}{}

\tagitem{prvOr}{%
  \indcase{!A}{(!B \lor !C)}{$\Struct{M} \VDash \indfrm$ precies als
    $\Struct{M} \VDash !B$ of $\Struct{M} \VDash !C$ (of beide).}}{}

\tagitem{prvIf}{%
  \indcase{!A}{(!B \lif !C)}{$\Struct{M} \VDash \indfrm$ precies als
    niet $\Struct {M} \VDash !B$ of $\Struct M \VDash !C$ (of beide).}}{}

\tagitem{prvIff}{%
  \indcase{!A}{(!B \liff !C)}{$\Struct{M} \VDash {\indfrm}$ precies
    als beide delen allebei gelden of allebei niet gelden. Dus: óf
    $\Struct{M} \VDash {!B}$ en $\Struct{M} \VDash {!C}$, óf noch
    $\Struct{M} \VDash {!B}$ noch $\Struct{M} \VDash {!C}$.}}{}

\tagitem{prvAll}{%
  \indcase{!A}{\lforall[x][!B]}{$\Struct{M} \VDash {\indfrm}$ precies
    als voor elk $a \in \Domain{M}$ geldt
    $\Struct{M[a/c]} \VDash \Subst{!B}{c}{x}$, waarbij $c$ niet
    in~$!B$ mag voorkomen.}}{}

\tagitem{prvEx}{%
  \indcase{!A}{\lexists[x][!B]}{$\Struct{M} \VDash
  {\indfrm}$ precies als er minstens één $a \in \Domain M$ is waarvoor
  $\Struct{M[a/c]} \VDash \Subst{!B}{c}{x}$, waarbij $c$ niet
  in~$!B$ mag voorkomen.}}{}
\end{enumerate}
Laat $x_1$, \dots, $x_n$ alle vrije !!{variable}s in~$!A$ zijn.
Kies nieuwe constantsymbolen $c_1$, \dots, $c_n$ die niet in~$!A$
staan. Kies ook $a_1$, \dots, $a_n \in \Domain M$ en laat
$s(x_i) = a_i$.

Bewijs dat $\Sat{M}{!A}[s]$ precies als
$\Struct M[a_1/c_1,\dots,a_n/c_n]
\VDash \Subst{\Subst{!A}{c_1}{x_1}\dots}{c_n}{x_n}$.

(Deze opgave laat zien dat je de betekenisregels voor eerste-ordelogica
ook zonder toekenningen aan variabelen kunt formuleren.)
\end{prob}

\begin{prob}
Stel dat $f$ een functiesymbool is dat nog niet in~$!A(x,y)$ staat.
Bewijs dat er !!a{structure}~$\Struct{M}$ bestaat waarvoor
$\Sat{M}{\lforall[x][\lexists[y][!A(x,y)]]}$, precies als er
!!a{structure}~$\Struct M'$ bestaat waarvoor
$\Sat{M'}{\lforall[x][!A(x,f(x))]}$.

(Dit is een bijzonder geval van de stelling van Skolem. De nieuwe
functie legt voor elke invoer een geschikte getuige vast. De formule
$\lforall[x][!A(x,f(x))]$ heet een
\emph{Skolemnormaalvorm} van
$\lforall[x][\lexists[y][!A(x,y)]]$.)
\end{prob}

\end{document}
